\chapter*{Daftar Simbol}

\begin{table}[htbp]
    \begin{tabular}{lcl}
         $R_t$      & : & \textit{Return} emas pada waktu ke-$t$ \\[4pt]
        $P_t$       & : & Harga emas pada waktu ke-$t$ \\[4pt]
        $P_{t-1}$   & : & Harga emas pada waktu sebelum $t$ \\[4pt]
        $r$         & : & Jangkauan data \\[4pt]
        $x_{min}$   & : & Nilai minimum \\[4pt]
        $x_{maks}$  & : & Nilai maksimum \\[4pt]
        $l$         & : & Lebar masing-masing kelas \\[4pt]
        $k$         & : & Jumlah kelas \\[4pt]
        $\Pr$       & : & Operator probabilitas \\[4pt]
        $X_n$       & : & Kondisi saat $t = n$ \\[4pt]
        $P_{ij}$    & : & Probabilitas transisi dari \textit{state} $i$ menuju \textit{state} $j$ \\[4pt]
        $n_{ij}$    & : & Jumlah transisi dari \textit{state} $i$ menuju \textit{state} $j$ \\[4pt]
        $N_i$       & : & Total transisi yang berasal dari \textit{state} $i$ \\[4pt]
        $S$         & : & Himpunan \textit{state} pada DTMC \\[4pt]
        $\mathbb{N}$ & : & Himpunan bilangan asli \\[4pt]
        $\square$   & : & Operator temporal \textit{always} \\[4pt]
        $\Diamond$  & : & Operator temporal \textit{eventually} \\[4pt]
        $s_0$       & : & \textit{State} awal pada model DTMC \\[4pt]
        $AP$        & : & Himpunan proposisi atomik \\[4pt]
        $L$         & : & Fungsi pelabelan \textit{state} dalam DTMC \\[4pt]
        $\Phi$      & : & Formula \textit{state} dalam logika temporal \\[4pt]
        $\phi$      & : & Formula \textit{path} dalam logika temporal \\[4pt]
        $U$         & : & Operator \textit{until} pada LTL/CTL/PCTL \\[4pt]
        $U_{\le n}$ & : & Batas waktu maksimal $n$ untuk operator hingga (\textit{until}) \\[4pt]
        $P_J(\phi)$  & : & Operator probabilitas PCTL dengan interval $J$ \\[4pt]
    \end{tabular}
\end{table}